\chapter{Inflation: The Evidence across Time and Countries}\label{ch:b3:inflation-evidence}

Inflation is the macroeconomic fact with the longest record and the clearest
regularities. This chapter sets the theory of the Year 1 and Year 2 volumes against
that record: prices under gold and under paper money, money growth and inflation
across countries, the hyperinflations and how they ended, the cost of bringing
inflation down, and what markets expect it to be. One thread runs through it:
lasting inflation is a monetary phenomenon whose cause is usually fiscal.

\begin{openingchart}{Consumer price level of five economies, 1870--2020 (1870 = 1,
logarithmic scale).}{Jordà, Schularick and Taylor, Macrohistory Database (cpi); our
rebasing}
\begin{tikzpicture}
\begin{semilogyaxis}[omtimesten, width=0.72\linewidth, height=0.46\linewidth, xmin=1865,
  xmax=2025, ymin=0.4, ymax=20000, ylabel={Price level (1870 = 1)}, legend pos=north west,
  legend style={font=\footnotesize}]
  \addplot[omDef, thick, mark=none] table[x=year, y=GBR] {figdata/bachelor-3/16-inflation-evidence/out/prices.csv};
  \addplot[black, thick, mark=none] table[x=year, y=USA] {figdata/bachelor-3/16-inflation-evidence/out/prices.csv};
  \addplot[omThm, thick, mark=none] table[x=year, y=FRA] {figdata/bachelor-3/16-inflation-evidence/out/prices.csv};
  \addplot[omThm, thick, dashed, mark=none] table[x=year, y=ITA] {figdata/bachelor-3/16-inflation-evidence/out/prices.csv};
  \addplot[omDef, thick, dashed, mark=none] table[x=year, y=JPN] {figdata/bachelor-3/16-inflation-evidence/out/prices.csv};
  \legend{United Kingdom, United States, France, Italy, Japan}
\end{semilogyaxis}
\end{tikzpicture}
\end{openingchart}

% figdata: prices_points
Under the \omterm{def:b3:exchange-rate-regimes:gold-standard}{gold standard}, prices in 1913 were about where they had been in 1870 ---
1.04 times the 1870 level in Britain, 0.76 times in the United States (Japan's, the
exception, doubled). Then they
rose in every economy: by 2020, 18 times the 1870 level in the United States and 89
times in Britain, and more than two thousand times in France, Italy and Japan,
whose price levels rose more than a hundredfold between 1913 and 1950. Four decades of stable prices, then a century of
rising ones: the difference is the monetary standard.

\begin{recall}
The Year 1 volume defined seigniorage, the inflation tax and hyperinflation, and
stated the quantity equation, the Fisher effect and the sacrifice ratio. The Year 2
volume treated rational expectations, credibility and the dynamics of government
debt. The grades 10--12 and Year 1 volumes told the American disinflation of
1979--1983; it is not retold here.
\end{recall}

\begin{notation}[Symbols of this chapter]\label{not:b3:inflation-evidence:symbols}
$\mu$ is the growth rate of money and $\pi$ the inflation rate, both continuous rates per
year; $m$ real money balances as a percent of GDP; $s$ seigniorage as a percent of
GDP; $d$ the primary deficit and $b$ real government debt; $\alpha$ the sensitivity of
money demand to expected inflation (a local symbol: not the capital share).
\end{notation}

\section{Money growth and inflation across countries}\label{sec:b3:inflation-evidence:quantity}

\begin{omchart}{Average inflation against average broad-money growth, 1990--2024
(percent a year, averages of log growth, logarithmic scales), 146 countries with
at least 20 years of data; the dashed line is equality. Grenada, whose money
supply shrank on average, cannot appear on a log axis.}{World Bank, World
Development Indicators (FM.LBL.BMNY.ZG, FP.CPI.TOTL.ZG); our computation}
\begin{tikzpicture}
\begin{loglogaxis}[omscatter, log ticks with fixed point, width=0.58\linewidth,
  height=0.46\linewidth, xmin=1, xmax=200, ymin=0.3, ymax=200,
  xlabel={Money growth (\% a year)}, ylabel={Inflation (\% a year)}]
  \addplot table[x=money, y=infl] {figdata/bachelor-3/16-inflation-evidence/out/mi.csv};
  \addplot[black, dashed, no marks] table[x=x, y=y] {figdata/bachelor-3/16-inflation-evidence/out/mi_line.csv};
\end{loglogaxis}
\end{tikzpicture}
\end{omchart}

% figdata: mi_stats
Across 146 countries, the regression of average inflation on average money growth,
both as continuously compounded rates, has a slope of 0.96 (standard error 0.02) and
an $R^2$ of 0.92. Most of the fit comes from the high-inflation countries; below 10~\%
a year, the points scatter and inflation runs well below money growth, because real
output grows and velocity moves. That is why, on the chart's log scales, the cloud
is steeper than the equality line: the log-log slope over the 145 countries shown is
1.36, with an $R^2$ of 0.83.

\begin{model}{The quantity theory over the long run}
\assumes Velocity changes slowly and output growth does not depend on money
growth in the long run.
\predicts Over decades, inflation equals money growth minus output growth plus the
change in velocity: one for one across countries with different money growth.
\tested Country averages over 35 years line up on a slope of 0.96, dominated by
high-inflation countries (the chart above).
\breaks Low-inflation samples, where velocity shifts and output growth dominate;
short periods, when money demand moves with interest rates and financial change.
\end{model}

\section{Hyperinflations}\label{sec:b3:inflation-evidence:hyper}

\begin{definition}[Cagan model]\label{def:b3:inflation-evidence:cagan}
\index{Cagan model}
In the \emph{Cagan model}, the demand for real money balances falls exponentially
with expected inflation: $m = Ae^{-\alpha\pi^e}$. In a steady state with $\pi^e = \pi = \mu$, the government
raises seigniorage $s = \mu m = \mu Ae^{-\alpha\mu}$.
\end{definition}

\begin{proposition}[The seigniorage curve]\label{prop:b3:inflation-evidence:laffer}
Steady-state seigniorage is maximised at $\mu = 1/\alpha$, where it equals $A/(\alpha e)$. Any revenue
below the maximum is raised by two rates of money growth.
\end{proposition}

\begin{proof}
$ds/d\mu = Ae^{-\alpha\mu}(1 - \alpha\mu)$, positive below $1/\alpha$ and negative above: $s$ rises, peaks at $1/\alpha$ with
$s = (1/\alpha)Ae^{-1}$, and falls toward zero as $\mu$ grows. A continuous function that rises
from 0 to its maximum and falls back toward 0 takes each intermediate value twice.
\end{proof}

\begin{omchart}{Steady-state seigniorage in the Cagan model with $A = 20~\%$ of GDP and
$\alpha = 2$ (the weekend problem's Kestria), against money growth.}{Our computation
(\code{cagan.py})}
\begin{tikzpicture}
\begin{axis}[omrecord, width=0.56\linewidth, height=0.36\linewidth, xmin=0, xmax=2.5,
  ymin=0, ymax=4.3, xlabel={Money growth $\mu$ (continuous, per year)},
  ylabel={Seigniorage (\% of GDP)}]
  \addplot[omDef, thick, mark=none] table[x=mu, y=s] {figdata/bachelor-3/16-inflation-evidence/out/laffer.csv};
  \draw[guide] (axis cs:0.5,0) -- (axis cs:0.5,3.68) node[above, font=\scriptsize] {peak at $1/\alpha$};
  \draw[guide] (axis cs:0,3) -- (axis cs:2.5,3) node[above left, font=\scriptsize] {a deficit of 3~\%};
\end{axis}
\end{tikzpicture}
\end{omchart}

The model explains why hyperinflations accelerate. A government whose deficit
exceeds the peak revenue cannot finance it at any steady rate; it prints faster,
people hold less money, and the base on which the tax falls shrinks. And it
explains how they end: when the deficit closes, money growth can fall to zero,
and the demand for real balances recovers.

\begin{record}[Germany, 1923]
In 1923 the German government financed virtually all of its spending by issuing
notes. On 15 October a decree created a new bank and currency, the Rentenmark,
worth one trillion paper marks, and capped both its total issue, at 3.2 billion,
and its lending to the government, at 1.2 billion; on 15 November the Reichsbank
stopped discounting treasury bills. In late November prices stopped rising and the
mark stopped falling, and the budget swung into balance. % ledger: b3ie:sar.germany
The same two measures ended the hyperinflations of Austria, Hungary and Poland: a
central bank legally bound to refuse the government new unsecured credit, and a
simultaneous change of fiscal regime. And in each case the note circulation went
on growing fast after prices had stabilised: people rebuilt the money balances
they had fled. % ledger: b3ie:sar.ends
\end{record}

The last sentence is the \omterm{def:b3:inflation-evidence:cagan}{Cagan model}'s remonetisation: once expected inflation
falls, real balances rise toward $A$, and the central bank can issue money without
inflation, once.

\section{The fiscal origin}\label{sec:b3:inflation-evidence:fiscal}

\begin{definition}[Fiscal theory of the price level; unpleasant arithmetic]\label{def:b3:inflation-evidence:fiscal}
\index{fiscal theory of the price level}\index{unpleasant arithmetic}
The \emph{fiscal theory of the price level} is the model in which the price level
adjusts so that the real value of government debt equals the present value of
future primary surpluses. The \emph{unpleasant arithmetic} is the result that,
when the fiscal authority does not adjust its deficits, tighter money now implies
higher inflation later.
\end{definition}

\begin{proposition}[Unpleasant arithmetic]\label{prop:b3:inflation-evidence:arithmetic}
A government with real debt $b_0$, a primary deficit $d$ in each of two periods and a
real rate $r$ must raise seigniorage $s_1$ and $s_2$ with $s_2 = (1 + r)\big((1 + r)b_0 + d - s_1\big) + d$ if
the debt is to be repaid at the end of period 2. Each unit less of seigniorage now
requires $1 + r$ units more later.
\end{proposition}

\begin{proof}
The debt at the end of period 1 is $b_1 = (1 + r)b_0 + d - s_1$. At the end of period 2 it must
be zero: $(1 + r)b_1 + d - s_2 = 0$. Substituting, $s_2 = (1 + r)b_1 + d$, and $\partial s_2/\partial s_1 = -(1 + r)$. With
seigniorage rising in money growth below the Cagan peak, the later money growth,
and so inflation, must be higher.
\end{proof}

This is the \omterm{def:b3:inflation-evidence:fiscal}{unpleasant arithmetic} of Sargent and Wallace (1981): a central bank
that fights inflation alone, against a government that does not close its deficit, only postpones
inflation and makes it larger. If money demand depends on expected inflation, even
the postponement fails: expected future inflation lowers the demand for money now
and raises inflation today.

\section{Disinflations and what they cost}\label{sec:b3:inflation-evidence:disinflation}

\begin{omchart}{Median inflation of the 26 inflation-targeting countries, by years from
the adoption of a target (year 0), and median inflation of all other countries in the
same calendar years (percent).}{Roger (2010) for adoption dates; World Bank, World
Development Indicators (FP.CPI.TOTL.ZG); our computation}
\begin{tikzpicture}
\begin{axis}[omtime, width=0.6\linewidth, height=0.38\linewidth, xmin=-5.5, xmax=8.5,
  ymin=0, ymax=12, xtick={-5,...,8}, xlabel={Years from adoption}, ylabel={Percent},
  legend pos=north east, legend style={font=\footnotesize}]
  \addplot[omDef, thick, mark=*] table[x=k, y=targeters] {figdata/bachelor-3/16-inflation-evidence/out/it_event.csv};
  \addplot[omThm, thick, dashed, mark=o] table[x=k, y=others] {figdata/bachelor-3/16-inflation-evidence/out/it_event.csv};
  \legend{Targeters, Other countries}
\end{axis}
\end{tikzpicture}
\end{omchart}

% figdata: it_event
Five years before adoption, the median targeter's inflation was 9.0~\%, against 6.0~\%
elsewhere in the same years; four years after, 2.8~\% against 3.9~\%. Inflation fell
everywhere over these decades, so the comparison is the gap: from 3.0 points above
the others to 1.1 below. Countries chose to target inflation when they had already
begun to bring it down, so the fall is not all the target's effect
(\cref{exo:b3:inflation-evidence:11}). % ledger: b3ie:roger.it

\begin{record}[Disinflations with and without credibility]
In 65 disinflations in moderate-inflation OECD economies, the output lost per point
of lasting inflation reduction --- the sacrifice ratio --- was smaller when
inflation was brought down fast than when it was brought down gradually, and
smaller where wages were set flexibly; the openness of the economy made no
difference. % ledger: b3ie:ball.sacrifice
The hyperinflations ended with almost no output cost, because the regime change was
believed; moderate inflations, whose end depends on contracts and expectations
adjusting, cost more.
\end{record}

\section{Measured expectations}\label{sec:b3:inflation-evidence:expectations}

\begin{definition}[Breakeven inflation]\label{def:b3:inflation-evidence:breakeven}
\index{breakeven inflation}
\emph{Breakeven inflation} is the yield on a nominal government bond minus the yield
on an inflation-indexed bond of the same maturity: the inflation rate at which the
two bonds pay the same.
\end{definition}

\begin{measure}[Breakeven inflation]
\omterm{def:b3:inflation-evidence:breakeven}{Breakeven inflation} equals expected inflation plus an inflation risk premium (what
holders of nominal bonds demand for bearing inflation risk) minus a liquidity
premium (indexed bonds trade less, so their yields are higher). In a panic, when
investors sell the less liquid indexed bonds, \omterm{def:b3:inflation-evidence:breakeven}{breakeven inflation} can collapse
without any change in expected inflation. Surveys of households and forecasters
measure expectations directly, with their own biases.
\end{measure}

\begin{omchart}{Ten-year breakeven inflation in the United States, monthly averages,
2003--2025 (percent): the nominal ten-year Treasury yield minus the ten-year
inflation-indexed yield.}{Board of Governors of the Federal Reserve System, H.15
(RIFLGFCY10\_N.M, RIFLGFCY10\_XII\_N.M); our difference}
\begin{tikzpicture}
\begin{axis}[omtime, width=0.7\linewidth, height=0.36\linewidth, xmin=2002, xtick distance=5, xmax=2026.2,
  ymin=0, ymax=3.2, ylabel={Percent}]
  \addplot[omDef, thick, mark=none] table[x=t, y=breakeven] {figdata/bachelor-3/16-inflation-evidence/out/breakeven.csv};
  \draw[guide] (axis cs:2003,2) -- (axis cs:2026,2);
\end{axis}
\end{tikzpicture}
\end{omchart}

% figdata: breakeven_stats
The breakeven rate averaged 2.06~\% from 2003 to 2019. It fell to 0.25~\% in December
2008, when the indexed market froze, and reached its highest, 2.89~\%, in April
2022, while consumer prices rose by 8~\% a year: ten-year expectations moved by less
than a point. Long-run expectations stayed anchored, which is the measured form of
the credibility the Year 2 volume modelled.

\section{Exercises}\label{sec:b3:inflation-evidence:exercises}

\begin{exercise}[$\star$]\label{exo:b3:inflation-evidence:1}
In the opening chart, identify a period of stable prices and the periods of fastest
rise for each economy. In what years did French and Italian prices rise most?
\end{exercise}

\begin{exercise}[$\star$]\label{exo:b3:inflation-evidence:2}
Prices rise by 50~\% a month. By how much do they rise in a year? In a day, on average?
\end{exercise}

\begin{exercise}[$\star$]\label{exo:b3:inflation-evidence:3}
The monetary base is 8~\% of GDP and grows by 25~\% a year with inflation of 22~\%. How
much seigniorage, as a share of GDP, does the government raise?
\end{exercise}

\begin{exercise}[$\star$]\label{exo:b3:inflation-evidence:4}
What episode would contradict the claim that hyperinflations have a fiscal
origin?
\end{exercise}

\begin{exercise}[$\star\star$]\label{exo:b3:inflation-evidence:5}
In the \omterm{def:b3:inflation-evidence:cagan}{Cagan model} with $\alpha = 2$ (years), find the revenue-maximising money growth as a
continuous rate and as an annual percentage. With $A = 15~\%$ of GDP, find the maximum
seigniorage.
\end{exercise}

\begin{exercise}[$\star\star$]\label{exo:b3:inflation-evidence:6}
With the JST data, compute decade averages of money growth and inflation for each
economy since 1870 and regress inflation on money growth. Compare the slope before
and after 1970.
\end{exercise}

\begin{exercise}[$\star\star$]\label{exo:b3:inflation-evidence:7}
In the two-period model, $b_0 = 50$, $d = 5$ and $r = 4~\%$. The central bank raises
$s_1 = 5$; find $s_2$. If it lowers $s_1$ to 2, by how much must $s_2$ rise?
\end{exercise}

\begin{exercise}[$\star\star$]\label{exo:b3:inflation-evidence:8}
Trend inflation falls from 9~\% to 3~\% over three years, during which output is 2, 3
and 1~\% below potential. Find the sacrifice ratio. A second episode lowers inflation
by the same amount in one year with a 4~\% output gap. Which was cheaper?
\end{exercise}

\begin{exercise}[$\star\star$]\label{exo:b3:inflation-evidence:9}
What evidence would show that \omterm{def:b3:inflation-evidence:breakeven}{breakeven inflation} is not expected inflation?
\end{exercise}

\begin{exercise}[$\star\star\star$]\label{exo:b3:inflation-evidence:10}
Prove the seigniorage proposition, and show that if expectations adapt slowly, a
government can raise more than the steady-state maximum for a while. Why can it not
do so for ever?
\end{exercise}

\begin{exercise}[$\star\star\star$]\label{exo:b3:inflation-evidence:11}
For the 26 targeters, compare average inflation in the five years before adoption
with the years 3 to 7 after, and do the same for all other countries in the same
calendar years. What is the difference-in-differences, and what would make it a
poor estimate of the target's effect?
\end{exercise}

\begin{exercise}[$\star\star\star$]\label{exo:b3:inflation-evidence:12}
A government ends a hyperinflation by closing its deficit and fixing its currency.
Before, real balances were 4~\% of GDP; with stable prices they return to 16~\% over
three years. How much money can the central bank issue without inflation? Why is it
dangerous to treat this as a lasting source of revenue?
\end{exercise}

\section{Problem: Kestria's Printing Press}\label{sec:b3:inflation-evidence:problem}

\begin{problem}[{Weekend problem --- Kestria's printing press}]\label{pb:b3:inflation-evidence:kestria}
Kestria's government runs a primary deficit of 3~\% of GDP and finances it by issuing
money. Real money balances are $m = 20e^{-2\pi}$ percent of GDP, with $\pi$ the continuous
inflation rate per year. % figdata: kestria_stats, laffer

\medskip\noindent\textbf{Part I --- the deficit and seigniorage.}
\begin{enumerate}
\item Write steady-state seigniorage as a function of money growth $\mu$.
\item Find seigniorage at $\mu = 0.1$ and $\mu = 1$.
\item Find the revenue-maximising money growth and the maximum seigniorage.
\item Express the revenue-maximising rate as an annual inflation rate.
\item Can Kestria finance a deficit of 4~\% of GDP by money in a steady state?
\item Why do real balances fall as inflation rises?
\end{enumerate}

\medskip\noindent\textbf{Part II --- the inflation that finances the deficit.}
\begin{enumerate}[resume]
\item Show that two rates of money growth raise 3~\% of GDP, and find them numerically.
\item Convert both into annual inflation rates.
\item Find real balances at each.
\item Which steady state is stable if expectations adapt slowly to actual
inflation? (Argue informally.)
\item Kestria is at the low steady state, and the deficit rises to 3.5~\% of GDP. What
happens to inflation?
\item What happens if the deficit rises above 3.68~\%?
\end{enumerate}

\medskip\noindent\textbf{Part III --- the stabilisation.} Kestria closes its deficit
and sets up a \omterm{def:b3:exchange-rate-regimes:regimes}{currency board}; inflation falls to zero and is expected to stay
there.
\begin{enumerate}[resume]
\item Find real balances after the reform.
\item Find the remonetisation: the money the board can issue without inflation, in
percent of GDP.
\item The board must back new money with reserves. Where can they come from?
\item Why is a \omterm{def:b3:exchange-rate-regimes:regimes}{currency board} alone not enough without the fiscal reform?
\item What did the German reform of 1923 have in common with this one?
\item Why did the note circulation keep growing after German prices stabilised?
\end{enumerate}

\medskip\noindent\textbf{Part IV --- a gradual alternative.} Instead of the reform,
Kestria brings inflation down gradually from the low steady state to 3~\% a year.
\begin{enumerate}[resume]
\item By how many points must inflation fall?
\item With a sacrifice ratio of 1.5, find the cumulative output loss.
\item With a credible reform and a sacrifice ratio of 0.5, find it.
\item What does the evidence on disinflations say about which ratio is likely?
\item During the gradual disinflation, how is the deficit financed?
\item Compare the output costs of the two strategies.
\item State the result: the inflation rate that finances Kestria's deficit, and the
remonetisation gain after the reform.
\end{enumerate}
\end{problem}
